\thispagestyle{empty}
\begin{center}
{\LARGE\bfseries 基于多源数据与结构水热耦合的\\[0.25em]吉隆泥石流触发及传播研究\par}
\vspace{0.65em}
{\large\bfseries 摘\quad 要}
\end{center}
\begingroup\small
针对高寒山区冰崩向泥石流演变的链式过程，本文建立“气象因素识别—结构水热稳定性—启动时间预测—传播响应”的多模型研究体系。

首先，基于MET、RADM、SOIL和FLUX构造5526个窗口、112项特征，通过边际异常、条件贡献与相似天气过程三路线识别因素族，得到土壤水分、空气湿状态、土壤热状态、空气热状态和平均风场的共识TOP5，并量化第五席竞争与排序稳定性。其次，在共同结构与参数设计下构建结构、结构加温度、结构加水及结构水热耦合四机制，完成2048个情景，形成条件临界关系；解析热核与细化有限差分的最大有限启动差为0.347小时，原始141条试验为本构竞争提供外部依据。

时间预测采用机制首次过阈进行条件状态回放，以逆速度/状态空间和动态生存模型在13个外部终末失稳组及1个真实右删失站上分组评价24、48、72小时预测。固定提前24小时发布72小时预测时，逆速度和状态空间方法的完整标签包含为3/13和6/13，显示信息尺度与发布设计的重要作用。传播部分完成34个简化和31个原生计算，以全源面ST、120米格网为主展示，比较孔压、复合固水、夹带及数值尺度响应；7个选定运行联接4个机械设计、49条目标记录。

研究以区域背景、条件机制、外部时间评价和给定释放传播分别回答四问，给出因素排序、临界关系、时间区间及过程场，并据此提出环境背景—源区状态—沟谷响应的分层监测配置。
\par\vspace{0.6em}\noindent\textbf{关键词：}因素族识别；结构水热耦合；首次过阈；时间预测；传播仿真
\par\vspace{0.7em}\noindent{\footnotesize AI辅助说明：本文部分文字整理、程序和图表使用OpenAI Codex辅助；具体贡献及使用记录见文后声明与配套说明。}
\endgroup
